\section{Inside an inference engine}
\label{sec:engines}
\subsection{From static batches to iteration-level scheduling}
In a static batch, requests can remain tied together until the longest member
finishes. A request producing ten tokens then wastes potential participation
while another produces a thousand. Iteration-level scheduling chooses the active
set again between model iterations, allowing finished requests to leave and
new work to enter. Orca established this design for generative-model
serving~\cite{orca}. ``Continuous batching'' is the common operational term.

Consider three requests with output lengths 2, 5, and 3. In a deliberately
simplified model where a three-slot iteration has constant cost, a fixed batch
runs five iterations for ten useful token updates out of fifteen possible
slots. A dynamic scheduler can fill vacated slots with new requests. This
example demonstrates available scheduling opportunity, not a $15/10$ speedup:
real iteration time changes with batch composition, prompt work, and cache
length. The scheduler also needs free state capacity before admitting work.

The engine's queue is not interchangeable with the external gateway queue.
Once admitted, a request may consume KV blocks and compete for token budget.
Holding a request outside the engine avoids some of that pressure but adds
waiting before execution. A useful gateway admission policy must understand
what the engine can already do, otherwise it can reduce batching and leave the
GPU idle.

\subsection{PagedAttention and the vLLM memory model}
The foundational vLLM design maps a sequence's logical KV blocks to physical
blocks allocated as needed. Blocks do not have to be adjacent in GPU memory.
This reduces large reservations and external fragmentation, and permits shared
blocks with reference counting and copy-on-write behavior~\cite{paged}.
Paging does not compress the tensors in Equation~\ref{eq:kvbytes}.

For a block capacity of $p$ tokens, sequence $i$ with $T_i$ tokens requires
$\lceil T_i/p\rceil$ blocks. Its unused tail is
\begin{equation}
 w_i=p\left\lceil T_i/p\right\rceil-T_i,\qquad 0\le w_i<p.
\end{equation}
Ignoring shared blocks, tail waste is less than $Bp$ token slots for $B$
sequences. Larger blocks reduce some metadata and indexing overheads but
coarsen allocation and reuse; smaller blocks can reduce unused tails while
increasing management work. These are design tradeoffs, not a universal best
block size. The appropriate choice also depends on the attention backend and
model's supported layouts.

\fig{papers/vllm}{1.0}{vLLM's original block-table translation diagram:
logical token blocks map to noncontiguous physical KV blocks. Reproduced from
Kwon et al.~\cite{paged}, Figure 6, arXiv:2309.06180v1, CC BY 4.0;
scaled and format-converted only.}{paging}

vLLM's versioned prefix-cache design identifies full blocks using a chain of
hashes over prior-prefix identity, current tokens, and relevant extra identity
fields. A new request can reuse a matching prefix; reference counts and a free
block structure distinguish currently used from evictable blocks~\cite{prefix}.
A prefix can be logically cached yet be evicted when space is needed. An active
reference protects state required by a current computation; retaining a
finished session's state is a separate policy problem.

\paragraph{Walkthrough: two related prompts.}
Suppose prompts are ``system + repository + task A'' and ``system + repository +
task B''. If tokenization and all computational identity match through the
repository portion, those prefix blocks can be reused. If the first differing
token is inside a block, a block-level implementation's reuse boundary may be
before that token. Hash lookup is not a semantic similarity search. Changing a
system instruction near the front can invalidate later exact-prefix reuse even
when most natural-language content looks similar.

\subsection{vLLM V1: token budgets and host overhead}
The V1 guide describes a unified scheduler assigning token counts per request,
rather than entirely separate scheduling abstractions for prompt and output
phases. This supports composing chunked prefill, prefix caching, and speculative
work within a token-budget model~\cite{v1}. The abstraction matters because the
engine can control work more finely than a gateway that only counts requests.
An 80-token prompt and an 80,000-token prompt are not equal demand.

A conceptual budget is
\begin{equation}
 \sum_{i\in\mathcal A}n_i^{\mathrm{scheduled}}\le K,
 \qquad M_{\mathrm{active}}+M_{\mathrm{reserved}}\le M_{\mathrm{KV,budget}}.
\end{equation}
A token count is only a proxy for work: decode at long context has a different
cost from a short-context token, and full, local, and latent attention differ.
The hardware ultimately constrains time and memory, not abstract token units.

The host still tokenizes, validates requests, maintains scheduling metadata,
prepares tensors, and processes outputs. A model server can use Python at the
orchestration layer while kernels and collectives execute in compiled code.
The relevant question is whether Python lies on the measured critical path.
Moving every component to Rust does not change the amount of KV data read by a
GPU or transferred over a link. Removing repeated scans, redundant RPCs, or
avoidable serialization can matter before a language change does.

\subsection{Chunked prefill and interference}
A large prefill inserted among active decodes can delay their next tokens.
Chunking divides prompt work across iterations, placing a bound on how much
prefill competes with decode at once. Sarathi-Serve develops scheduling around
chunked prefills to manage this throughput/latency tension~\cite{sarathi}.

Let a chunk have $c$ tokens and estimated execution time $t_p(c,B,T)$ under the
current decode set. A conceptual ITL budget imposes
\begin{equation}
 t_d(B,T)+t_p(c,B,T)+t_{\mathrm{other}}\le\tau_{\mathrm{ITL}}.
\end{equation}
This additive approximation is conservative when work overlaps and optimistic
when it ignores interference. Smaller $c$ creates more scheduling boundaries
and may hurt compute efficiency. Larger $c$ can improve prompt throughput while
worsening decode tail gaps. Adaptive chunking requires measured response surfaces,
not simply choosing the smallest chunk.

\subsection{Other engine designs and why comparison matters}
\textbf{SGLang.} Its original work combines structured LLM-program execution
with RadixAttention, representing reusable prefixes in a radix-tree-oriented
cache~\cite{sglang}. This is an alternative organization of prefix reuse, not
an absence of memory constraints. Its current HiCache design extends reuse
beyond GPU memory; Section~\ref{sec:hierarchy} separates that layer from the
original engine mechanism.

\textbf{TensorRT-LLM.} The documented KV subsystem includes block pools, reuse,
secondary-memory offload, and prioritized eviction. Retention can attach
priorities to token ranges~\cite{trt}. Therefore a blanket statement that
``engines expose no retention controls'' is incorrect. ATFM must discover which
control is available in the deployed backend and compare against that engine's
best native retention policy.

\textbf{vAttention.} The memory-management paper with this name keeps the
virtual address layout contiguous while mapping physical GPU memory on demand
using CUDA virtual-memory APIs~\cite{vattention}. It addresses fragmentation
without requiring the same non-contiguous virtual layout as PagedAttention.
This paper is distinct from a later sparse-attention project with the same
name. The broader lesson is that block paging is one way to manage growing
state, not a mathematical prerequisite for efficient inference.

\subsection{Speculation, quantization, and parallel execution}
Speculative decoding generates candidate tokens with a cheaper draft process
and verifies them using the target model. With the appropriate acceptance and
correction algorithm it can preserve the target sampling distribution~\cite{spec}.
A simple throughput comparison is
\begin{equation}
 t_{\mathrm{effective\ token}}\approx
 \frac{t_{\mathrm{draft}}(k)+t_{\mathrm{verify}}(k)+t_{\mathrm{control}}}
      {\E[A(k)]},
\end{equation}
where $A(k)$ counts emitted tokens per round under the exact implementation,
including any additional verified token. Acceptance, draft overhead, available
batching, and memory use determine whether the method helps. A request-level
capacity estimator must not assume one target forward pass always equals one
emitted token.

Weight quantization reduces model storage/traffic; KV quantization reduces
sequence-state traffic. They alter different terms in the budget. KIVI studies
asymmetric low-bit quantization of keys and values~\cite{kivi}; low-bit cache
storage requires numerical-quality evaluation, scale metadata, and suitable
kernels. A factor-of-eight reduction in scalar bit width does not imply an
end-to-end factor-of-eight memory or latency reduction.

Parallelism redistributes work rather than removing it. Tensor parallelism
splits operations and adds collectives; pipeline parallelism splits layers and
can suffer bubbles; data parallelism creates independent model replicas and
needs routing; context parallelism spreads long sequence work; expert parallelism
routes mixture-of-experts tokens among expert shards. A fleet may combine them.
A cache-local routing decision that ignores collective topology can be worse
than a small cache miss on a better-connected worker group.

\takeaway{Before adding a higher-level controller, establish a tuned engine
baseline with continuous batching, appropriate kernels, realistic token budgets,
and the native cache/reuse features supported by the chosen model. ATFM should
be evaluated on top of that baseline, not credited for optimizations already
available in it.}
